\section{Related Work}
\label{sec:related}

\subsection{Parameter-Space Model Merging and Task Arithmetic}
The task of consolidating multiple task-specific expert neural networks into a single multi-task model has gained substantial traction due to its computational efficiency. Initial investigations centered around \emph{Model Soups}~\cite{wortsman2022model}, which demonstrated that averaging the weights of multiple models fine-tuned from the same pre-trained initialization on the same task improves accuracy and out-of-distribution generalization. This was extended to disparate tasks by \emph{Task Arithmetic}~\cite{ilharco2022editing}, where fine-tuned weights are subtracted from the pre-trained base model to form "task vectors" that can be linearly added back to the base model. To resolve the negative interference or "parameter masking" that occurs when multiple task vectors are combined, TIES-Merge~\cite{yadav2023ties} introduced sign consensus filtering and parameter pruning. Other works have explored coordinate-alignment techniques such as ZipIt!~\cite{stoica2023zipit} or optimal transport barycenters~\cite{matena2022merging} to align representations before merging. While these approaches have shown strong empirical promise, they rely heavily on static, global heuristic ensembling coefficients that do not account for the high heterogeneity across different network layers.

\subsection{Few-Shot Adaptation and the Overfitting-Optimizer Paradox}
To address layer-wise parameter sensitivity, recent methods have proposed optimizing layer-wise or parameter-wise merging coefficients. Under extreme few-shot or test-time conditions, these coefficients are optimized using a small calibration dataset $\mathcal{D}_{\text{cal}}$ or an unlabeled test-time data stream~\cite{trial2_submission1, trial3_submission2}. However, as highlighted by \emph{PolyMerge}~\cite{trial2_submission3} and \emph{RegCalMerge}~\cite{trial2_submission1}, directly optimizing high-dimensional coefficient spaces on small calibration batches exposes the model to the \emph{Overfitting-Optimizer Paradox}. When the optimization search space is overparameterized relative to the extremely small sample size, the optimizer overfits to the transductive sampling noise of the calibration set. Although the training or calibration loss is minimized, the resulting ensembling parameters exhibit chaotic inter-layer oscillations, which catastrophically degrades the model's out-of-distribution generalization capability.

\subsection{PAC-Bayesian Generalization Theory}
The PAC-Bayesian framework, originally introduced by McAllester~\cite{mcallester1999pac} and further developed by Catoni~\cite{catoni2007pac} and others~\cite{langford2001pac, alquier2021user}, provides a rigorous learning-theoretic framework for bounding the generalization gap of randomized classifiers. By specifying a prior distribution $P$ over the hypothesis space before observing the training data, and optimizing a posterior distribution $Q$ based on the data, PAC-Bayesian bounds express the generalization risk in terms of the empirical training risk and the Kullback-Leibler (KL) divergence between $Q$ and $P$. Historically, this theory has been applied to bound the generalization of deep neural networks by injecting Gaussian noise into the network weights, demonstrating that "flat" regions of the loss landscape possess superior generalization guarantees~\cite{dziugaite2017computing, neyshabur2017exploring, guedj2019primer}.

While previous work has applied learning-theoretic concepts to model merging---most notably Rademacher-Bounded Polynomial Merging (RBPM)~\cite{trial5_submission2} which employed an empirical Rademacher complexity bound---their choice of regularization has remained largely heuristic. Specifically, RBPM's Consensus-Pulling $L_1$ regularizer was chosen intuitively and corresponds to a Laplace-like prior over the polynomial coordinates. In contrast, \textbf{PAC-Bayes Merge} provides the first formal, information-theoretic justification for trajectory-regularized model merging, mathematically proving that a Gaussian PAC-Bayesian prior centered at the uniform ensembling consensus analytically justifies a smooth quadratic $L_2$ Consensus-Pulling penalty.

\subsection{Stochastic Weight Averaging and Flat Minima}
Our work is also deeply connected to the literature on flat minima and Stochastic Weight Averaging (SWA)~\cite{izmailov2018averaging}. SWA proves that averaging multiple points along the SGD trajectory leads to a wider, flatter basin of attraction, which significantly improves out-of-distribution generalization. This principle underlies many domain generalization methods, such as SWAD~\cite{cha2021swad}, and sharpness-aware optimization techniques like SAM~\cite{foret2021sharpness}. In our work, we establish a formal connection between uniform weight merging and SWA, showing that the high performance of the Static Uniform baseline under extreme data noise is a direct consequence of the low-pass parametric filtering effect of weight averaging, which cancels out random independent task-specific optimization noise.
